\begin{center}
  \textsc{Abstract}
\end{center}
\noindent

Hospitals accumulate large volumes of biomedical documents that record diagnoses, treatments, test results, and the reasoning that connects them. Structuring this information would support clinical research, but the text is highly sensitive, rarely shareable, and almost absent from public French resources. This thesis investigates how compact information-extraction models can be trained without hospital records and deployed within the institutions where those records remain.

We first determine which public text can provide relevant training signal. We assemble and release French biomedical corpora from public sources. Because their documents mix relevant medical content with boilerplate, we annotate them at paragraph level by content type and quality. Selecting the relevant paragraphs matches full-corpus pretraining with one third of the tokens and yields two million public clinical-case passages. We also show that generic neural quality filters can amplify benchmark contamination.

We then train compact biomedical encoders. Continuing a general French model's pretraining improves all five biomedical benchmarks tested without training a new one from scratch. We next introduce a causal detour strategy: the model learns from next-token prediction, then returns to standard bidirectional masked language modelling. This beats matched MLM training on all eight French tasks and transfers to English. We also introduce OntoBook, which complements text pretraining by converting medical ontologies into synthetic textbooks; it improves document classification and medical coding. We release the resulting models and data.

To turn these encoders into useful systems, we address the downstream task that motivates this thesis: extracting study-specific variables from clinical documents. Each project asks for different fields, usually with little or no annotated data. We use large language models to annotate public text, to judge unseen types when fixed references fall short, and to generate synthetic lymphoma dossiers derived from clinical cases published in scientific articles. Clinicians read these dossiers through a web-based annotation tool and found them plausible overall. The generated supervision trains compact open-vocabulary extractors whose target-field descriptions are supplied as text at inference. We release the general-purpose extractors and a synthetic benchmark, on which our smallest task-adapted model outperforms extractive QA and comes within 0.018 value-$F_1$ of the strongest fine-tuned generator, a model twenty-seven times its size.

Together, these results show that compact clinical extractors can be built from carefully selected public data and synthetic supervision, without treating model scale as a substitute for data design. Transfer to real hospital records remains to be fully established. Finally, we explore how the rapid evolution of foundation models, including the emergence of agents that act across records and tools, makes realistic evaluation, especially of clinical safety, an urgent research priority.

\clearpage
\selectlanguage{french}
\begin{center}
  \textsc{Résumé}
\end{center}
\noindent

Les hôpitaux accumulent des documents biomédicaux qui consignent les diagnostics, les traitements, les résultats d'examens et le raisonnement qui les relie. Structurer ces informations faciliterait la recherche clinique, mais ces textes sont très sensibles, rarement partageables et presque absents des ressources publiques en français. Cette thèse étudie comment entraîner des modèles compacts d'extraction d'information sans dossiers hospitaliers, puis les déployer là où ces dossiers demeurent.

Nous déterminons d'abord quels textes publics fournissent un signal d'entraînement pertinent, et publions des corpus biomédicaux français. Comme leurs documents mêlent contenu médical et contenu générique, nous les annotons par paragraphe selon le type et la qualité. Sélectionner les paragraphes pertinents égale le préentraînement complet avec un tiers des tokens et fournit deux millions de passages publics de cas cliniques. Nous montrons aussi que les filtres neuronaux génériques de qualité peuvent amplifier la contamination des jeux d'évaluation.

Nous entraînons ensuite des encodeurs biomédicaux compacts : poursuivre le préentraînement d'un modèle généraliste français améliore les cinq jeux d'évaluation biomédicaux étudiés, sans repartir de zéro. Nous introduisons un détour causal : le modèle apprend d'abord par prédiction du prochain token, puis revient au masquage bidirectionnel standard. Il surpasse un entraînement MLM comparable sur les huit tâches françaises et se transfère à l'anglais. Nous introduisons aussi OntoBook, qui convertit des ontologies médicales en manuels synthétiques ; il améliore la classification de documents et le codage médical. Nous publions les modèles et les données obtenus.

Pour en faire des systèmes utiles, nous abordons la tâche qui motive cette thèse : extraire de documents cliniques des variables propres à chaque étude, avec des champs différents à chaque projet et peu de données annotées. Nous utilisons de grands modèles de langue pour annoter des textes publics, juger des types jamais observés faute de références fixes et produire des dossiers synthétiques de lymphome à partir de cas cliniques publiés. Des cliniciens les ont consultés dans un outil en ligne et jugés globalement plausibles. Cette supervision entraîne des extracteurs compacts à vocabulaire ouvert, dont les champs cibles sont décrits en texte à l'inférence. Nous publions ces extracteurs et un jeu d'évaluation synthétique, sur lequel notre plus petit modèle adapté à la tâche surpasse les questions-réponses extractives et se situe à 0,018 de value-$F_1$ du meilleur générateur affiné, vingt-sept fois plus grand.

Ces résultats montrent que des extracteurs cliniques compacts se construisent à partir de données publiques sélectionnées et d'une supervision synthétique, sans faire de la taille du modèle un substitut à la conception des données. Leur transfert aux dossiers hospitaliers réels reste à établir. Enfin, nous explorons comment l'évolution rapide des modèles de fondation, notamment des agents agissant sur des dossiers et des outils, fait de l'évaluation réaliste, en particulier de la sécurité clinique, une priorité urgente.
\selectlanguage{english}
